\section{A support test and the boundary of identifiability}
The portable six-phase Multi-RPM fixture supplies 218 measured current/speed
locations. Evaluate the resistance sensitivity $\rotationgenerator\current/
\omegae$ at those points. A nuisance voltage error $-a\current$ produces
the identical column, so the stacked two-column design has rank one.
Adding two constant dq voltage-offset columns gives rank three of four;
the exact confounder remains. This result is generated from the fixture,
not inferred from the appearance of its measured flux residuals.
Two speed slices cannot separate two identically scaled columns.

The synthetic tests use a known magnetic reference. On real data, angle
columns depend on flux and differential inductance that require independent
qualification. Current errors change both coordinates and values; an apparent
local correction can reflect an erroneous reference derivative. Rank after
whitening and parameter scaling matters more than visual separation.

Inverter dead time, nonlinear device drops, timing, losses, thermal state and
genuine asymmetry can add nuisance directions \cite{liu2018,richter2014}.
They are not all identified by this model. Hysteresis and dynamic magnetic
states invalidate a static current-only reference. Full rank in a restricted
model gives local conditional identifiability, not global or causal truth.
The measured support test validates an operator property while preserving
that distinction.
